% Section 10: selection framework (P4)
The stack of \cref{sec:stack} tells you \emph{what a system is}, and this section tells a
practitioner \emph{what to build}. The selection framework routes by task; model preference is not
an axis. It is written as a
procedure because the assembly is the product. Each of the six steps hands the next a term it
cannot invent for itself, and the last one ends at a configuration whose residual risk is named
and owed.

\subsection{The two axes}
\label{sec:sel-axes}
The horizontal axis is \textbf{answer-space openness}: how far the correct answer set is from
enumerable at design time. Well-defined answer spaces support closed contracts (typed
decisions, validation); open-ended ones cannot be closed by any contract layer and must rely
on grounding and review (\cref{sec:case-rag}). We use \emph{openness of the
answer space}. Stochastic inference is a deployment knob; an axis built on it would re-ask the
determinism question and leave every reliability limit unaddressed.

The vertical axis is \textbf{cost of error}: the severity of a wrong output. Reversible low-stakes
outputs allow sampling or iteration strategies; irreversible high-stakes ones push the operating
point toward validated acceptance and human escalation, whichever objective family serves the
task. \cref{fig:map} draws the resulting quadrants and the escalation band.

\subsection{The procedure}
\label{sec:sel-procedure}
The procedure starts from the paper's one falsifiable cross-layer regularity: failures originate
where they are least catchable and are usually caught in the later columns of
\cref{fig:failmatrix}. Where the checks sit is therefore a decision the procedure makes. It
consumes the ledger.

\begin{enumerate}[leftmargin=1.6em, itemsep=2pt]
\item \textbf{Risk specification.} Fix the three terms that make ``reliable'' a claim in this
  paper: the envelope of failure kinds the system must not produce, the budget it must hold, and
  the assurance path that intervenes as it drifts (\cref{sec:reliability}). Mapped onto the
  framework, those terms are two axes---answer-space openness and error severity---plus two
  constraints, reversibility and operational budget: three claims, which the framework reads as
  four terms every later step reads.

\item \textbf{Placement.} Locate the task on \cref{fig:map}: how open the answer set is at design
  time, and what a wrong output costs. The coordinate is an analytic cut [J], the status the map
  gives its quadrant lists. Record the reason the placement holds rather than treating the
  quadrant as measured, and keep task openness separate from decoding stochasticity, which is a
  deployment knob (\cref{sec:reliability}).

\item \textbf{Modifier routes.} Walk the four placement modifiers --- distribution shift,
  reversibility, observability, operational budget --- plus the schedule route, which asks the
  question the axes never raise: \emph{which factorization does this contract have to inspect?}
  The schedule is not inherited from the objective; \cref{fig:schedule-atlas} draws the three
  factorizations as lanes, and \cref{sec:sel-modifiers}(v) states the route. Each route is stated
  with its evidence there; any of them can move the operating point without
  moving the axes, so the placement from Step~2 is provisional until they have been walked.

\item \textbf{Probe against the failure matrix.} For each failure kind in the envelope, find the
  row where it originates in \cref{fig:failmatrix} and ask which later column would catch it;
  where no column does, the candidate needs another layer or a narrower claim, and the residue is
  bracketed as the residual of \cref{fig:layer5-assur}. A candidate that verifies only at the
  layer where it operates is checking the wrong place by default: origins cluster at L2 while
  filled cells cluster in the last three columns. The two probes already on
  \cref{fig:map} are our own cases, and they are illustrative: they show what a filled placement
  looks like before any threshold exists.

\item \textbf{Archetype routes.} Read the quadrant's candidate modes and select one to carry:
  (low cost, well-defined) rules, classifiers and decision models benchmarked against baselines;
  (high cost, well-defined) specialist plus validators, calibrated acceptance and escalation;
  (high cost, open-ended) grounded generative assistance with retrieval, verification and human
  review; (low cost, open-ended) general-purpose generation with sampling and iteration. Two gated
  variants enter only where their precondition holds: latent-prediction planning under
  identifiability conditions (\cref{sec:case-jepa}), and rules or formal verification where an
  executable validator exists (\cref{sec:case-forms}). Then read the card of
  \cref{sec:familycards} --- what the objective optimizes, which task types it suits, which
  failure modes remain, which complementary layers cover the residual --- and the mechanism
  sketch for where the design's benefit comes from. The route fixes the objective family and the
  contract tier; it does not fix the schedule of Step~3.

\item \textbf{Budget covenant.} Convert the L0 budget into thresholds and escalation conditions
  with the existing [T]-class machinery --- conformal abstention \cite{id240501563},
  risk-controlled selective answering \cite{id260704430}, and risk-budgeted serving
  \cite{id260933887} --- then fit them on the operator's calibration set, reference configuration
  and correctness measure. The cutoff is the cost-based reject option, chosen to hold the error
  budget on the calibration distribution \cite{geifman2017selective}; never install it as a fixed
  universal constant (see also the contested-direction caveat \cite{id260933843}); the output is an
  \emph{operator}, not a ``safe zone''. Instrument it: monitors, abstention and escalation ladder,
  plus a scheduled evaluation plan that tracks shift and takeover attempts, the contested claims of
  \cref{tab:crosscase} included.
\end{enumerate}
What the six steps produce is a candidate operating mode with a checked residual-risk story.
The framework draws two lines: it treats ranking model classes against each other as a
category error (\cref{sec:stack}), and it holds escalation to the residual the evidence supports.

\subsection{Modifier routes}
\label{sec:sel-modifiers}
Four modifiers move the placement and leave the axes fixed. A fifth route, the schedule, does a
different job: it selects which factorization the contract inspects. All five are claims we have
to defend, and each answers a question the placement alone cannot settle.

\begin{description}[leftmargin=1.5em, itemsep=1pt]
\item[(i) Distribution shift.] Widens margins on the reliability side: a well-instrumented
  classifier degrades unpredictably under shift, making calibrated abstention the floor rather
  than a bonus \cite{id250315850,id221116158}. Shift is also the matrix's one ambiguous origin row,
  filed at L0 and stated as such (\cref{fig:failmatrix}).
\item[(ii) Reversibility.] Converts some high-error-cost placements back into the iterative
  regime, because repairable damage can wait for repair. The
  dashed elbow of \cref{fig:map} is this move drawn.
\item[(iii) Observability.] Determines whether the assurance layer gets signal at all; monitoring
  under low observability is one of the honest failed rows of \cref{sec:case-classifier}, where
  the monitor is evaded under ordinary task pressure \cite{id260930217}.
\item[(iv) Operational budget.] Constrains the acceptable cost of the safeguard chain itself; it
  constrains which configuration is even selectable, not how reliable it should be. A
  configuration that reaches target quality faster converts its margin into reliability budget
  (more candidates, more checks) rather than raw throughput; the schedule classes of
  \cref{fig:schedule-atlas} make that exchange rate visible, and on XSum the NCM authors spend
  saved sampling time on minimum-Bayes-risk selection for \evaurep{about 7\% higher ROUGE-L within
  the same sampling budget} (\cref{sec:l5-selfcheck}).
\item[(v) Schedule.] The route that puts the schedule axis into the procedure: it fixes which
  factorization the contract must inspect, and so where the acceptance surface sits. A schedule is
  a factorization choice, not an architecture: the same objective admits a left-to-right chain, an
  iterative denoising sweep, or a coarse-to-fine level order, and what differs across the three is
  how many sequential decisions a contract must accept and whether any of them can be taken back
  (\cref{fig:schedule-atlas}, \cref{sec:contracts}). The archetype of Step~5 therefore
  fixes the objective family; the schedule is a separate commitment, and it is
  the one that decides whether the budget conversion of (iv) buys more passes or more checks. Where
  a shorter path holds quality, the saved passes are spendable on verification: the hierarchical
  family reports 21$\times$ faster decoding on LM1B at matched parameter scale and 31$\times$ on
  XSum at matched ROUGE-L, at $K \ll L$ [A] \cite{kumbong2026ncm}, and its tokenizer-stage
  substitutions are what the contract's exactness boundary then has to catch.
\end{description}

\subsection{A worked instance: typed decision serving}
\label{sec:sel-instance}
The six steps are easier to check when one task is carried through them. This instance carries
Case~2 (\cref{sec:case-decision}), a non-generative decision model returning
\texttt{choice}/\texttt{score}/\texttt{noul} from one forward pass. It adds no measurement of its
own; the single cell the paper cannot supply is marked operator-supplied.

\begin{description}[leftmargin=1.5em, itemsep=2pt]
\item[\textbf{1. Risk specification.}] A closed decision over an option set of 50 similar
  candidates, the size at which the inherited evaluation reports its accuracy drop; a 10\%
  accepted-error budget; no human reader in the loop, because the decision acts. Two failure kinds
  fill the envelope: interface instability under option order and option-set size, and acceptance
  miscalibrated against the budget.
\item[\textbf{2. Placement.}] High cost of error against a well-defined answer space: the quadrant
  routing to a specialist with validators, calibrated acceptance and escalation (\cref{fig:map}).
  The placement is an analytic cut \textsf{[J]}: it records that a wrong decision acts and that
  the option set is enumerable at design time.
\item[\textbf{3. Modifier routes.}] Shift is presentational here: a paired evaluation holds the
  inputs fixed, reverses option order, and 30\% of answers flip, with accuracy dropping at 50
  similar candidates \evindep{\cite{id261002267}}. Reversibility is none once the decision acts, so
  the gate sits before commitment. Observability is the serving harness's request and decision
  logs. Budget is one forward pass, leaving no sampling passes to convert into checks. The schedule
  route does not arise: a single-pass decision has one commitment point.
\item[\textbf{4. Probe against the failure matrix.}] Interface instability is the decision-model
  framing row: it originates at L4 and is caught in the contracts column, where shuffling against
  order and a shortlist against set size both sit. The row's monitoring-and-escalation cell is an
  empty gap in our evidence, so the gate this instance adds is a design move; no catch is
  documented. Acceptance
  miscalibration originates at L1, and its acceptance cell is filled only with a
  contested direction. The joint effect of the three is unmeasured, and no anchor evaluates
  harness, gate and escalation together (\cref{sec:case-decision}): the procedure brackets that
  gap, it does not close it.
\item[\textbf{5. Archetype route.}] Within that route the family is the decision-oriented one, at
  the typed-decision contract tier, with the card of \cref{sec:familycards} supplying the failure
  modes it leaves and the complements that cover them. The validators cover form. The
  contract settles that the answer is a well-formed choice over the offered set and relocates
  whether the choice is correct.
\item[\textbf{6. Budget covenant.}] The published gate freezes its thresholds before evaluation and
  still misses its 10\% accepted-error target on both reported tracks, though it beats random
  escalation \evindep{\cite{id260933843}}. The shipped checkpoint is independently measured
  under-confident, mean confidence 0.553 against accuracy 0.767, signed gap $-0.214$
  \evindep{\cite{id260933843}}. That sign contradicts the card's over-confidence framing and decides
  which way a confidence-gated cascade fails, so the operator refits the cut on its own calibration
  split \evvendor{a threshold is a policy chosen from measured accuracy on the operator's own data,
  not a property of the model}. Coverage at the 10\% budget is not reported in this paper; that cell
  is operator-supplied. Escalation: any accepted set that misses the budget routes to human review,
  with a measured catch rate.
  Owner: the operator. Re-evaluation trigger: any checkpoint change or execution-configuration
  change, or any option set above 50.
\end{description}
Shuffling, a shortlist and a recalibrated gate have not been composed end to end. That residue is
why the instance stops short of a guarantee and hands back a configuration with a named gap and an
owner for it.

\subsection{Why this is useful beyond standard practice}
Existing deployed reliability practice solves \emph{per-dimension} problems
\cite{id250206872,id240910102,id230901219}: a definition each for calibration, RAG honesty, or content
quality; each works, is checkable, has failure modes; assembling them for a concrete task is
what breaks. Our framing explicitly conditions on L0 and treats the assembly as the product:
each safeguard layer declares what it covers, what its unverified residues are
(\cref{fig:failmatrix}), and which verified substitutes exist if a residue is deal-breaking.
The matrix narrows the candidates to the verification-anchored set it enumerates, each with its
ledger record, and the evidence map (\cref{fig:evidence}) tells the reader which of these claims
are contested.

The regularity can be tested outside our own matrix. We coded four incidents chosen for their
public record rather than their fit---Amazon's recruiting tool (2018), the Dutch childcare-benefits
risk model (2013--2021), Zillow Offers (2021), and the Air Canada chatbot (2024)---with the coding
rule fixed before the cases were assigned: origin is the earliest layer whose documented design
choice made the harm likely, and a catch is a safeguard inside the deployed system that fired
before harm~\cite{reuters2018amazon,amnesty2021xenophobic,zillow2021winddown,aircanada2024bccrt}.
In all four the earliest documented origin is L1 or L2, and none was caught at or before that
layer: three had no internal catch at all, and the fourth was stopped only by project cancellation
at L5. Read at the strength the sample supports, origins cluster early; the sample is too small to
locate the mode within L1--L2. The later columns describe safeguards that exist rather than
safeguards that fire. Zillow's documented L4 co-origin
is the boundary case, and counting nominal review steps as catches would inflate the catch
side.\footnote{One coder, coding rule fixed before assignment, not blinded to the matrix; a
retrospective coding, not a benchmark.} The coding manual and every assignment are
\cref{app:codebook} (\cref{tab:codebook}).
